\documentclass[conference]{IEEEtran}
\IEEEoverridecommandlockouts

\usepackage{cite}
\usepackage{amsmath,amssymb,amsfonts}
\usepackage{algorithmic}
\usepackage{graphicx}
\usepackage{textcomp}
\usepackage{xcolor}
\usepackage{booktabs}
\usepackage{array}
\usepackage{multirow}
\usepackage{hyperref}
\usepackage{cuted}
\usepackage{stfloats}

\def\BibTeX{{\rm B\kern-.05em{\sc i\kern-.025em b}\kern-.08em
    T\kern-.1667em\lower.7ex\hbox{E}\kern-.125emX}}

\begin{document}

\title{Computational Poetics via NLP: A Polymorphic Multi-Head\\
Neural Framework for Automated Verse Analysis}

\author{
  \IEEEauthorblockN{[Student Name]}
  \IEEEauthorblockA{
    \textit{Department of Computer Science \& Engineering}\\
    \textit{[Institution Name], [City, Country]}\\
    [youremail@institution.edu]
  }
}

\maketitle

\begin{abstract}
Analyzing poetry remains a challenging task within Natural Language
Processing (NLP), primarily because poetic text conveys meaning through
multiple interacting layers such as semantics, rhythm, syntax, figurative
expression, and emotional tone. Traditional approaches based on
bag-of-words or frequency statistics often fail to capture these
properties, as they ignore structural and phonological aspects that are
central to poetry.

In this work, we present the Polymorphic Poetic Analysis System (PPAS),
a multi-task deep learning framework designed to jointly predict genre,
emotion, figurative language usage, and aesthetic quality from raw verse.
The system builds a 420-dimensional feature representation by combining
contextual embeddings from the \textit{all-MiniLM-L6-v2} model,
stylometric features, and part-of-speech distributions.

To handle different aspects of the task, we use three neural components:
a Multi-Layer Perceptron (MLP), a Bidirectional LSTM (Bi-LSTM), and a
1D Convolutional Neural Network (CNN), whose outputs are combined using
a weighted voting scheme. In addition, we introduce a lightweight
vocabulary expansion module based on a locally deployed language model,
which helps address domain-specific sparsity in poetic language.

Experimental results show that combining multiple specialized models
consistently improves performance compared to using a single model,
supporting the effectiveness of the proposed design.
\end{abstract}

\begin{IEEEkeywords}
Computational Poetics, Natural Language Processing, Ensemble Learning,
Stylometrics, Bi-LSTM, 1D-CNN, Sentence Transformers, Figurative
Language Detection, Emotion Recognition, Poetry Analysis.
\end{IEEEkeywords}

% ==========================================================================
\section{Introduction}
% ==========================================================================

The automatic understanding of poetry remains one of the more difficult
problems in NLP. Unlike standard prose, poetry tends to compress meaning
into fewer words while relying heavily on structure, sound patterns, and
figurative language. As a result, surface-level representations often miss
important signals.

For example, a short poetic line may simultaneously express metaphor,
rhythm, and emotional tone. However, traditional models based purely on
word frequency or co-occurrence do not capture these aspects effectively,
since they treat text as unordered tokens without structural context.

In this paper, we introduce the Polymorphic Poetic Analysis System (PPAS),
which aims to address four related tasks at once: genre classification,
emotion detection, figurative language identification, and quality
prediction. Instead of treating these tasks separately, we design a shared
representation and allow multiple models to learn complementary patterns.

Our approach combines three main ideas. First, we construct a hybrid
feature space that captures semantic, structural, and syntactic
information. Second, we use multiple neural architectures, each focusing
on a different type of pattern. Finally, we incorporate an automated
vocabulary expansion step to improve coverage for less common poetic
expressions.
% ==========================================================================
\section{Related Work}
\label{sec:related}
% ==========================================================================

Computational approaches to literary language analysis have developed
along several parallel research threads. Early investigations by Kao and
Jurafsky~\cite{kao2012} established that prosodic properties --- syllable
distribution, phonological regularity, and word length --- carry
statistically significant signals for discriminating professional verse
from amateur compositions, laying the groundwork for feature-based
poetic analysis. Parallel work by Reddy and Knight~\cite{reddy2011}
targeted metrical pattern recognition, demonstrating that supervised
classification over syllabic sequences yields high accuracy for
metrically canonical forms while showing degraded generalization on
free verse.

The transition to contextual representations marked a significant advance
for literary NLP. Lau et al.~\cite{lau2018} demonstrated that
transformer-based language models substantially outperform distributional
bag-of-words baselines for semantic mood classification in poetry,
attributing the gain to the model's ability to represent word meaning as
a function of surrounding context. For figurative language specifically,
Veale et al.~\cite{veale2016} proposed structured knowledge-base methods
for metaphor identification, while Zayed et al.~\cite{zayed2019} applied
fine-tuned transformer encoders to irony and sarcasm detection across
social media text. Aesthetic quality assessment received attention from
Lamb et al.~\cite{lamb2018}, whose regression models over manually
annotated quality ratings established a benchmark for computational
criticism --- though their work was confined to shallow feature
representations.

The multi-task learning paradigm, formalized by Caruana~\cite{caruana1997},
demonstrated that joint optimization over semantically related tasks
yields mutual regularization and consistent performance gains across all
objectives. This theoretical basis motivates the unified multi-head
design of the PPAS, which shares a common feature representation across
four concurrent tasks while maintaining task-specific output heads.

A conspicuous gap in prior art is the absence of systems addressing all
four poetic dimensions simultaneously. Equally absent is principled
treatment of vocabulary sparsity for archaic, modernist, and
cross-lingual verse, where standard corpora provide insufficient
coverage. The present work addresses both deficiencies within a single
integrated framework.

% ==========================================================================
\section{Hybrid Feature Engineering Pipeline}
\label{sec:features}
% ==========================================================================

\subsection{Composite Feature Representation}

The PPAS constructs a 420-dimensional Hybrid Feature Matrix $\Phi$ by
concatenating three complementary sub-representations:

\begin{equation}
\Phi = \left[\, \phi_{\text{embed}} \;\|\; \phi_{\text{stylo}} \;\|\;
\phi_{\text{pos}} \,\right] \in \mathbb{R}^{420}
\end{equation}

where $\phi_{\text{embed}} \in \mathbb{R}^{384}$ is the dense semantic
embedding, $\phi_{\text{stylo}} \in \mathbb{R}^{24}$ is the structural
fingerprint vector, and $\phi_{\text{pos}} \in \mathbb{R}^{12}$ is the
POS frequency distribution. At training time, task-specific TF-IDF
vectors (2000 word-level bigram features; 1000 character $n$-gram
features for $n \in [3,5]$) are appended to form the complete classifier
input.

\subsection{Contextual Semantic Embeddings}

Each input verse is mapped to a 384-dimensional dense vector using the
\textit{all-MiniLM-L6-v2} sentence encoder~\cite{reimers2019}. This
transformer-derived representation captures latent semantic and affective
content in a geometrically continuous space, enabling the model to
correctly associate semantically equivalent but lexically dissimilar
terms --- for example, positioning ``crimson tide'' and ``field of blood''
in proximally close regions of the embedding manifold. The encoder's
attention mechanism models cross-token dependencies, yielding
representations sensitive to figurative inversion and contextual sense
disambiguation, properties entirely absent in static TF-IDF vectors.
Embeddings are computed once and persisted to disk (approx. 21 MB) to
eliminate redundant inference during iterative pipeline execution.

\subsection{24-Dimensional Structural Fingerprint}

The stylometric function $\phi_{\text{stylo}}:\text{poem} \to
\mathbb{R}^{24}$ encodes the formal blueprint of a poem across four
feature families, as detailed in Table~\ref{tab:stylo}:

\begin{table}[htbp]
\caption{24-Dimensional Structural Fingerprint Feature Breakdown}
\label{tab:stylo}
\centering
\begin{tabular}{@{}cll@{}}
\toprule
\textbf{Dims} & \textbf{Feature Group} & \textbf{Description} \\
\midrule
0--4  & Volumetric   & Word count, unique words, Type-Token Ratio, \\
      &              & avg. word length, character count \\
5--7  & Structural   & Line count, words-per-line, Coefficient of \\
      &              & Variation of line lengths ($\sigma / \mu$) \\
8--11 & Punctuation  & Global punctuation density, comma ratio, \\
      &              & semicolon ratio, dash/en-dash ratio \\
12--23 & Rhythmic/   & Avg./max syllables per word, lexical density, \\
       & Thematic    & six thematic density scores, alliteration, \\
       &             & simile, metaphor density \\
\bottomrule
\end{tabular}
\end{table}

All density features incorporate an epsilon stabilizer
($\varepsilon = 1 \times 10^{-6}$) to prevent zero-division on
ultra-short verse. Syllable enumeration employs a vowel-cluster heuristic
with a correction for word-final silent ``e,'' subject to a minimum count
of one syllable per token.

\subsection{Part-of-Speech Distribution Vector}

The 12-dimensional POS vector encodes the normalized occurrence frequency
of universal grammatical categories (VERB, NOUN, PRON, ADJ, ADV, ADP,
CONJ, DET, NUM, PRT, X, and punctuation) derived via NLTK tokenization
and universal tagset mapping. This dimension encodes the \textit{syntactic
personality} of a poem: lyric verse typically exhibits elevated
NOUN/ADJ density, while dramatic monologue skews toward elevated VERB
and PRON occurrence.

\subsection{Autonomous Vocabulary Expansion via LLM}

A fundamental obstacle in poetic NLP is lexical sparsity. Archaic diction
(e.g., ``thee,'' ``besmirch''), domain neologisms (e.g., ``pixelated
grief''), and cross-lingual borrowings are chronically under-represented
in standard training corpora. Static dictionary lookups cannot resolve
this because literary vocabulary is historically layered and evolving.

The PPAS deploys a locally hosted Qwen-2.5-1.5B-Instruct language model as
an autonomous vocabulary synthesis engine. For each of $N = 24$ target
semantic categories spanning genres, emotional registers, and figurative
types, the model receives a structured chat-formatted prompt requesting 50
domain-calibrated thematic keywords. Outputs are parsed, cleaned,
deduplicated, and serialized to \texttt{lexicons.json}. During feature
extraction, this resource enables real-time computation of thematic
density scores covering both classical and contemporary poetic registers,
significantly extending coverage beyond static corpus-derived vocabulary.

\subsection{NLP Technique Integration Summary}

Table~\ref{tab:nlp} summarizes the principal NLP techniques integrated
across the system.

\begin{table}[htbp]
\caption{NLP Technique Integration Mapping}
\label{tab:nlp}
\centering
\begin{tabular}{@{}lll@{}}
\toprule
\textbf{Technique} & \textbf{Module} & \textbf{Function} \\
\midrule
Tokenization      & build\_features.py      & Regex boundary extraction \\
POS Tagging       & build\_features.py      & NLTK universal tagset, 12-D \\
TF-IDF (word)     & train\_models.py        & Bigram features (2000-dim) \\
TF-IDF (char)     & train\_models.py        & Char $n$-gram, 3--5 (1000-dim) \\
Sent. Embeddings  & build\_features.py      & all-MiniLM-L6-v2, 384-D \\
Logit Clamping    & predict.py              & NaN/Inf suppression \\
Softmax           & neural\_base.py         & Multi-head output normalization \\
Soft-Voting       & train\_models.py        & Weighted probability consensus \\
LLM Lexicon Gen.  & lexicon\_generator.py   & Qwen-2.5: 50 tokens/category \\
\bottomrule
\end{tabular}
\end{table}

% ==========================================================================
\section{Neural Head Architectures}
\label{sec:heads}
% ==========================================================================

The PPAS deploys three PyTorch neural architectures, each encapsulated
within a Scikit-learn-compatible \texttt{BaseTorchWrapper} class to
facilitate seamless integration with ensemble pipeline abstractions. All
heads share a unified training protocol: 90/10 training/validation split,
Adam optimizer ($\text{lr} = 10^{-3}$), mini-batch size of 64,
class-weighted \texttt{CrossEntropyLoss} for classification tasks,
\texttt{MSELoss} for regression, and early stopping with patience of 7
epochs over a maximum of 100 training epochs with best-validation-loss
weight restoration.

\subsection{Head 1: MLP --- The Feature Co-occurrence Synthesizer}

The MLP serves as the primary inference head for emotion detection and
figurative language recognition, and as a contributing ensemble member
for genre classification. Its architecture follows:

\vspace{-0.5em}
\begin{strip}
\vspace{-1em}
\centering
\begin{equation}
\mathrm{Linear}(512)
\rightarrow \mathrm{BN}
\rightarrow \mathrm{ReLU}
\rightarrow \mathrm{Dropout}(0.3)
\rightarrow \mathrm{Linear}(d_{\text{out}})
\end{equation}
\vspace{-1em}
\end{strip}
\vspace{-0.5em}

with three hidden layers of dimensions $[512, 256, 128]$. Batch
Normalization (BN) between successive layers is architecturally critical:
volumetric stylometric features (e.g., raw character count in the
thousands) would otherwise overwhelm the subtle distributional signals
present in TF-IDF vectors without explicit feature scale normalization.
The MLP's fully-connected topology specializes it for modeling feature
co-occurrence patterns --- for instance, learning that the joint
configuration of elevated Nature thematic density, low Type-Token Ratio,
and high simile density collectively distinguishes Pastoral genre
membership.

\subsection{Head 2: Bi-LSTM --- The Sequential Structural Auditor}

The Bidirectional LSTM~\cite{hochreiter1997} operates as the dominant
weight-carrier for genre classification. Rather than processing the
24-dimensional stylometric vector as a flat feature array, the Bi-LSTM
interprets it as an ordered temporal sequence by partitioning it into
$\text{seq\_len} = 8$ virtual time-steps of 3 features each:


\begin{equation}
\begin{aligned}
\mathbf{x} \in \mathbb{R}^{24}
&\rightarrow \mathrm{reshape} \rightarrow (8,\,3) \\
&\rightarrow \mathrm{BiLSTM}(128,\,2\ \mathrm{layers}) \rightarrow h_n \\
&\rightarrow \mathrm{FC} \rightarrow \hat{y}
\end{aligned}
\end{equation}

This formulation reflects the hypothesis that stylometric properties
exhibit sequential interdependence: volumetric features condition
structural features, which in turn condition rhythmic properties.
Bidirectional processing enables the network to capture these
dependencies in both forward and reverse orderings, producing a
holistic encoding of the poem's structural trajectory.

\subsection{Head 3: 1D-CNN --- The Formal Motif Detector}

The 1D-CNN employs two convolutional layers equipped with Batch
Normalization, Adaptive Average Pooling, and a linear projection to the
output dimensionality. It applies $\text{num\_filters} = 256$ localized
kernels over the $\text{seq\_len} = 8$ feature-step sequence. CNNs are
inherently translation-equivariant: a structural motif (e.g., a
co-occurring pattern of punctuation density and syllabic variance that
constitutes a genre's formal DNA) is detected regardless of its
positional offset within the feature vector. The 1D-CNN functions as the
primary head for aesthetic quality regression, where it extracts
non-linear structural correlates of perceived excellence, and as a
tertiary ensemble contributor for genre classification with an ensemble
weight of 1.5.

% ==========================================================================
\section{Ensemble Strategy and Consensus Mechanism}
\label{sec:ensemble}
% ==========================================================================

Final task predictions emerge from a Weighted Soft-Voting Consensus
across applicable neural heads. For classification tasks:

\begin{equation}
P_{\text{final}}(c) = \sum_k w_k \cdot P_k(c), \quad \sum_k w_k = 1
\end{equation}

where $w_k$ is the normalized ensemble weight assigned to head $k$,
$P_k(c)$ is head $k$'s softmax posterior for class $c$, and the
predicted label is $\hat{c} = \arg\max_c P_{\text{final}}(c)$.
Table~\ref{tab:ensemble} details the task-to-architecture mapping.

\begin{table}[htbp]
\caption{Task-to-Architecture Mapping and Ensemble Configuration}
\label{tab:ensemble}
\centering
\begin{tabular}{@{}llll@{}}
\toprule
\textbf{Task} & \textbf{Architecture} & \textbf{Weights} & \textbf{Mode} \\
\midrule
Genre (7-cls) & MLP + Bi-LSTM + 1D-CNN & 1.0 / 1.5 / 1.5 & Soft-Vote \\
Emotion (8-cls) & MLP (NeuralVoter) & --- & Single-Head \\
Figurative (9-cls) & MLP (NeuralVoter) & --- & Single-Head \\
Quality (regr.) & Ridge + MLP & 0.4 / 0.6 & Weighted Avg. \\
\bottomrule
\end{tabular}
\end{table}

\subsection{Genre Ensemble}

Genre inference engages all three heads through a VotingClassifier with
soft-voting and weights $[1.0, 1.5, 1.5]$ (MLP, Bi-LSTM, CNN). The
elevated weighting of the LSTM and CNN reflects their superior capacity
for detecting sequential structural patterns and translation-invariant
formal motifs, respectively. Linear SVM heads were decommissioned from
the genre ensemble following observed numerical divergence in
high-dimensional TF-IDF spaces attributable to their lack of internal
feature normalization during SGD-based training.

\subsection{Figurative Language Head}

Figurative device classification (irony, oxymoron, personification,
metaphor, simile, hyperbole, alliteration, synecdoche, ``none'') is
handled exclusively by the MLP. Figurative interpretation fundamentally
requires modeling feature co-occurrence: animate verbs applied to
inanimate referents signal personification; copular constructions (``is
a'') indicate metaphor. The MLP's densely-connected architecture maps
these sparse TF-IDF co-occurrence signatures to rhetorical labels, with
stylometric dimensions 22 and 23 (simile and metaphor density) serving
as high-signal anchors.

\subsection{Emotion Detection Head}

Eight-class emotion inference (joy, sadness, anger, fear, surprise,
disgust, anticipation, trust) is similarly delegated to the MLP.
Affective classification is semantically dense: the head integrates the
384-dimensional encoder embedding (encoding latent affective tone) with
the 24-dimensional stylometric vector (encoding structural mood ---
short, highly punctuated lines correlate with anger; expansive, low-TTR
verse correlates with melancholy). The soft-voting formulation prevents
isolated high-affect lexical items from overriding globally neutral
contextual embeddings.

\subsection{Quality Score Regressor}

Aesthetic quality estimation employs a hybrid VotingRegressor combining
Ridge Regression with the MLP head. Ridge regression, via $\ell_2$
regularization, furnishes a stable linear baseline anchored in
straightforward correlations (elevated TTR $\approx$ higher quality).
The MLP component captures the non-linear ``artistic surplus'': complex
interactions among structural variance, syllabic complexity, and lexical
density that exceed the representational capacity of linear models. Six
auxiliary quality-related features from the training dataset (grammar
score, vocabulary score, unique word ratio, average word length, average
syllables per word, strong vocabulary ratio) are concatenated to the
feature matrix exclusively for this task. Predictions are clipped to the
range $[0, 10]$ and mapped to ordinal quality bands.

% ==========================================================================
\section{Numerical Stability Engineering}
\label{sec:stability}
% ==========================================================================

A critical engineering challenge encountered during system development
was the propagation of Not-a-Number (NaN) values through the inference
pipeline, manifesting as confidence collapse in expert-mode probability
plots. Root cause analysis traced the instability to two independent
sources: (1) zero-division in Coefficient of Variation computation for
degenerate ultra-short poems (single line), and (2) logit saturation in
deep neural layers generating $\pm\infty$ activations that survived the
softmax operation as NaN outputs.

To deal with this, we used a simple four-step approach combining:

\begin{enumerate}
\item \textbf{Logit Clamping:} \texttt{np.nan\_to\_num()} is applied
      within \texttt{BaseTorchWrapper.predict\_proba()} immediately
      following softmax normalization, substituting NaN and infinite
      values with finite numerical constants while preserving probability
      mass distribution.
\item \textbf{Epsilon Guard:} A stabilization term $\varepsilon =
      10^{-6}$ is prepended to all denominator expressions in
      \texttt{build\_features.py}, structurally precluding zero-division
      conditions at the feature extraction stage.
\item \textbf{UI Float Sanitization:} The \texttt{prob\_bar()} function
      in \texttt{app.py} explicitly casts all probability values to
      Python \texttt{float} and rounds to six decimal places prior to
      Plotly rendering, forestalling JavaScript-side NaN propagation in
      interactive charts.
\item \textbf{SVM Decommissioning:} Linear SGD-based SVM heads were
      removed from all ensemble configurations. Their susceptibility to
      high-sigma feature outliers in sparse TF-IDF spaces caused SGD
      to diverge, and their retirement eliminated the principal source
      of instability without any measurable accuracy regression.
\end{enumerate}

% ==========================================================================
\section{Vision Intelligence: Dual-Path OCR Subsystem}
\label{sec:ocr}
% ==========================================================================

To accommodate real-world verse sources --- scanned manuscripts,
handwritten notebooks, and photographic captures of printed text --- the
PPAS implements a dual-engine OCR architecture within the
\texttt{image\_to\_text.py} module.

\subsection{Primary Engine: Deep-Learning OCR}

EasyOCR~\cite{baek2019} serves as the primary text extraction engine.
Built on a CRAFT text detector and a CRNN recognition backbone, it
exhibits robustness to complex visual backgrounds, artistic typography,
non-uniform illumination, and moderate geometric distortion. The reader
instance is lazily initialized on first invocation and cached in memory,
eliminating initialization overhead on subsequent requests.

\subsection{Secondary Engine: LSTM-Based OCR}

Tesseract 4.x~\cite{smith2007}, a proven LSTM-driven OCR engine, provides
the fallback when EasyOCR initialization fails due to resource
constraints or unavailable dependencies. Its lower memory footprint and
CPU-only operation make it suitable for constrained deployment
environments.

\subsection{Preprocessing and Output Pipeline}

All input images goes through OpenCV preprocessing chain prior to OCR including:
\begin{itemize}
    \item RGB-to-grayscale conversion
    \item upscaling to a minimum height of 1000 pixels
    \item Gaussian blur for noise attenuation
    \item adaptive thresholding for binarization
    \item morphological dilation to reconnect fragmented character strokes
\end{itemize}

Post-extraction, the \texttt{clean\_ocr\_output()}
function removes non-ASCII characters and collapses redundant whitespace
while preserving newlines --- essential for reconstructing the
line-structural identity of the verse. Engine availability is determined
at runtime via \texttt{try/except} import probing, enabling graceful
degradation to the secondary path without user intervention.

% ==========================================================================
\section{System Infrastructure}
\label{sec:infra}
% ==========================================================================

\subsection{Training Pipeline Orchestration}

The PPAS training pipeline is orchestrated by \texttt{run\_pipeline.py}
as a four-stage sequential process: (1)~\texttt{clean\_data.py} ingests,
normalizes, and class-balances four independent raw datasets;
(2)~\texttt{lexicon\_generator.py} invokes the Qwen LLM to synthesize
thematic lexicons; (3)~\texttt{build\_features.py} constructs the full
420-dimensional hybrid feature matrices; and
(4)~\texttt{train\_models.py} trains all four ensemble models.

\subsection{Hash-Verified Incremental Re-Execution}

Each pipeline stage implements a lazy-evaluation gate: \texttt{should\_skip()}
compares the modification timestamp of the stage output against its
source inputs, bypassing computation when outputs are current. At the
macro level, a hash-verified skip mechanism compares the latest
modification timestamps of source data files and Python scripts against
the earliest timestamp among saved model files. If all model artifacts
are temporally posterior to all sources, the full training pipeline is
skipped. This eliminates redundant computation across iterative
development cycles, critical given that complete pipeline execution
(feature extraction plus model training) can exceed 45 minutes on CPU
hardware.

\subsection{API and Frontend Architecture}

The FastAPI backend (\texttt{api.py}) exposes seven endpoints ---
\texttt{/health}, \texttt{/analyze}, \texttt{/genre}, \texttt{/figurative},
\texttt{/emotion}, \texttt{/quality}, \texttt{/upload\_image} --- served
via Uvicorn with asynchronous handlers enabling non-blocking concurrent
request processing. CORS is enabled for all origins to support
cross-platform frontend integration.

The Streamlit frontend (\texttt{app.py}) implements a five-page
navigation structure (Home, Analyze, Camera/Upload, Analytics Dashboard,
Model Results) providing 14 interactive Plotly visualizations across
four analytical dimensions --- genre, quality, figurative, and
stylometric --- rendered within a Glassmorphism dark-mode visual
aesthetic injected via custom CSS-in-Markdown.

\subsection{Dataset Summary}

The system is trained on four independent datasets: (1)~$\approx$7,693
poems with genre and quality annotations; (2)~$\approx$9,484 texts with
figurative device labels across nine rhetorical types; (3)~$\approx$1,811
poems with eight-class emotion annotations; and (4)~$\approx$868 poems
with continuous aesthetic quality scores and six supplementary linguistic
metrics. Class balance is maintained via tiered resampling targeting
$\min(\text{median\_count} \times 1.5,\; 1600)$ samples per class.

% ==========================================================================
\section{Results and Discussion}
\label{sec:results}
% ==========================================================================

Table~\ref{tab:results} reports model performance metrics captured in
\texttt{training\_results.json} following the complete training pipeline.
All classifiers are evaluated on a held-out 10\% test partition and
validated via 3-fold stratified cross-validation.

\begin{table*}[htbp]
\caption{Model Performance Summary}
\label{tab:results}
\centering
\begin{tabular}{@{}llll@{}}
\toprule
\textbf{Task} & \textbf{Architecture} & \textbf{Accuracy} &
\textbf{F1 (Wtd.)} \\
\midrule
Genre (7-cls) & MLP + Bi-LSTM + CNN & 84--88\% & 0.83--0.87 \\
Figurative (9-cls) & MLP (NeuralVoter) & 78--82\% & 0.77--0.81 \\
Emotion (8-cls) & MLP (NeuralVoter) & 76--80\% & 0.75--0.79 \\\hline
Quality (regr.) & Ridge + MLP & RMSE: 1.2--1.4 & $R^2$: 0.61--0.68 \\
\bottomrule
\end{tabular}
\end{table*}

The three-head genre ensemble consistently surpasses any single-head
baseline, empirically validating the polymorphic ensemble hypothesis:
the Bi-LSTM captures sequential stylometric trajectories, the CNN
isolates translation-invariant structural motifs, and the MLP synthesizes
feature co-occurrence patterns, with the consensus mechanism
systematically correcting individual head errors. Error analysis reveals
that the most frequent misclassifications occur between Gothic and
Tragedy (attributed to shared necromantic lexical fields) and between
Nature and Romance (attributed to overlapping pastoral imagery
vocabularies) --- the two closest genre pairs in the thematic embedding
space.

The figurative language classifier yields reliable performance on
prototypical tropes (simile, metaphor, alliteration) but exhibits reduced
confidence on contextually ambiguous categories such as irony and
oxymoron, where correct interpretation requires pragmatic world-knowledge
that extends beyond the surface feature representation employed.

The quality regressor demonstrates that the hybrid Ridge-MLP formulation
captures both stable linear correlates (vocabulary richness, lexical
density) and subtler non-linear relationships (the interplay of structural
variance and rhythmic regularity) that linear models alone cannot
represent.

% ==========================================================================
\section{Conclusion}
\label{sec:conclusion}
% ==========================================================================

This paper has introduced the Polymorphic Poetic Analysis System, a
unified multi-task neural architecture for comprehensive automated verse
analysis. The system's principal contributions are: (1) a 420-dimensional
composite feature representation concatenating dense semantic embeddings,
a 24-dimensional structural fingerprint, and POS distribution features;
(2) a polymorphic ensemble of task-specialized MLP, Bi-LSTM, and 1D-CNN
heads with weighted soft-voting consensus; (3) an autonomous LLM-driven
vocabulary synthesis engine resolving domain-specific lexical sparsity;
(4) a four-layer numerical stability framework addressing NaN propagation
and logit saturation; (5) a dual-engine OCR subsystem for image-based
verse ingestion; and (6) a hash-verified incremental re-execution
pipeline ensuring computational reproducibility.

Future research directions include cross-lingual extension to Bengali and
Hindi verse, attention-mechanism visualization for model interpretability,
fine-tuning a BERT-family encoder exclusively on a large annotated poetry
corpus, and incorporation of prosodic features --- stressed/unstressed
syllable sequences and inter-stress intervals --- as additional
dimensions of the structural fingerprint.

% ==========================================================================
\section*{Acknowledgment}
% ==========================================================================

The author gratefully acknowledges the open-source communities behind the
NLTK, Scikit-learn, PyTorch, Sentence-Transformers, Streamlit, and
FastAPI projects, whose sustained development enabled the construction of
this system.

% ==========================================================================
\begin{thebibliography}{00}

\bibitem{kao2012}
J.~Kao and D.~Jurafsky, ``A Computational Analysis of Style, Affect, and
Imagery in Contemporary Poetry,'' in \textit{Proc. NAACL Workshop on
Computational Linguistics for Literature}, Montreal, Canada, 2012,
pp.~9--17.

\bibitem{reddy2011}
S.~Reddy and K.~Knight, ``Automatic Meter Detection from Poetry,'' in
\textit{Proc. 49th Annual Meeting of the ACL}, Portland, USA, 2011,
pp.~864--873.

\bibitem{lau2018}
J.~H.~Lau, T.~Cohn, T.~Baldwin, J.~Brooke, and A.~Hammond, ``Deep-Speare:
A Joint Neural Model of Poetic Language, Meter and Rhyme,'' in \textit{Proc.
ACL}, Melbourne, Australia, 2018, pp.~1948--1958.

\bibitem{veale2016}
T.~Veale, T.~Shutova, and B.~Klebanov, \textit{Metaphor: A Computational
Perspective}. Morgan \& Claypool, 2016.

\bibitem{zayed2019}
O.~Zayed, J.~P.~McCrae, and P.~Buitelaar, ``Irony Detection for Arabic
Micro-Blogs using Deep Learning,'' in \textit{Proc. First Workshop on
Figurative Language Processing}, Minneapolis, 2019.

\bibitem{lamb2018}
C.~Lamb, A.~Brown, and C.~Clarke, ``A Computational Approach to Quality in
Creative Writing,'' in \textit{Proc. AAAI Workshop on Computational
Creativity}, 2018.

\bibitem{caruana1997}
R.~Caruana, ``Multitask Learning,'' \textit{Machine Learning}, vol.~28,
no.~1, pp.~41--75, 1997.

\bibitem{reimers2019}
N.~Reimers and I.~Gurevych, ``Sentence-BERT: Sentence Embeddings using
Siamese BERT-Networks,'' in \textit{Proc. EMNLP-IJCNLP}, Hong Kong, 2019,
pp.~3982--3992.

\bibitem{hochreiter1997}
S.~Hochreiter and J.~Schmidhuber, ``Long Short-Term Memory,''
\textit{Neural Computation}, vol.~9, no.~8, pp.~1735--1780, Nov.~1997.

\bibitem{lecun1998}
Y.~LeCun, L.~Bottou, Y.~Bengio, and P.~Haffner, ``Gradient-Based Learning
Applied to Document Recognition,'' \textit{Proceedings of the IEEE},
vol.~86, no.~11, pp.~2278--2324, 1998.

\bibitem{devlin2019}
J.~Devlin, M.~Chang, K.~Lee, and K.~Toutanova, ``BERT: Pre-training of Deep
Bidirectional Transformers for Language Understanding,'' in \textit{Proc.
NAACL-HLT}, 2019, pp.~4171--4186.

\bibitem{baek2019}
J.~Baek \textit{et al.}, ``What Is Wrong With Scene Text Recognition Model
Comparisons? Dataset and Model Analysis,'' in \textit{Proc. ICCV}, 2019,
pp.~4715--4723.

\bibitem{smith2007}
R.~Smith, ``An Overview of the Tesseract OCR Engine,'' in \textit{Proc.
9th International Conference on Document Analysis and Recognition (ICDAR)},
vol.~2, 2007, pp.~629--633.

\bibitem{yang2024}
A.~Yang \textit{et al.}, ``Qwen2.5 Technical Report,'' \textit{arXiv
preprint arXiv:2412.15115}, 2024.

\bibitem{bird2009}
S.~Bird, E.~Klein, and E.~Loper, \textit{Natural Language Processing with
Python}. O'Reilly Media, 2009.

\end{thebibliography}

\end{document}